\section{Continuous Proof, CI, and Test Explosion}

Agent-native testing should be organized by proof question, not author convenience. Domain invariants need Rust unit and property tests. Application workflows need authorization, idempotency, and transaction tests. Adapters need integration and contract tests. PostgreSQL needs migration, constraint, and schema drift checks. The browser needs component tests, rendered UX proof, and Playwright coverage for critical paths. Public benchmarks and open-source agents show that setup reliability, context retrieval, and interface design are central bottlenecks for autonomous repair \cite{sweAgentAci2024,sweBench2024,setupBench2025,contextBench2026,openhands2024,aiderGithub2026}.

\begin{table}[t]
\caption{Minimum proof lanes by ownership cell.}
\label{tab:proof-lanes}
\centering
\small
\begin{tabularx}{\columnwidth}{L{0.28\columnwidth} Y}
\toprule
\textbf{Layer} & \textbf{Minimum proof} \\
\midrule
Rust domain & Unit tests plus property tests for invariants and state machines. \\
Rust application & Integration tests for authz, idempotency, transactions, workflows. \\
Rust adapters & DB/external integration tests and failure injection. \\
Contracts & Generation check, drift check, compatibility check. \\
PostgreSQL & Migration apply, constraint tests, RLS policy tests where used. \\
TypeScript web & Typecheck, component tests, rendered UX QA, Playwright critical paths. \\
Python AI service & Contract tests, eval suites, no direct product-truth ownership. \\
Ops/security & Secret scan, dependency scan, provenance, audit gate. \\
\bottomrule
\end{tabularx}
\end{table}

The explosion of tests is useful only if routed. A repo with thousands of tests and no \texttt{test-map.json} still asks agents to guess. The standard therefore requires a fast lane for changed files, a contract lane for generated interfaces, a security lane for dependency and secret risk, a DB lane for migrations, a web lane for component and rendered UX proof, an e2e lane for critical browser behavior, and a release lane for full confidence.

CI modes should be explicit. \emph{Advisory} mode emits reports but does not fail merge. \emph{Guarded} mode blocks critical caps. \emph{Standard} mode enforces high and critical findings plus a score floor. \emph{Ratchet} mode prevents score regression without a named exception. \emph{Release} mode gates shipped artifacts on audit, contracts, security, DB, e2e, and version checks.

The continuous audit loop is concrete and local. A changed path enters \texttt{agent/owner-map.json} for ownership and \texttt{agent/test-map.json} for lane routing. \texttt{jankurai lane} or \texttt{jankurai proof} turns that route into a proof plan. \texttt{jankurai prove} executes allowed commands, writes command receipts and \texttt{target/jankurai/evidence-index.json}, and records artifact digests. \texttt{jankurai proof-verify} checks the plan and evidence against the current repository state. \texttt{jankurai audit} then folds ownership, generated zones, proof receipts, security evidence, UX evidence, and score history into JSON/Markdown/SARIF/JUnit reports. \texttt{jankurai witness} compares the candidate against an accepted baseline and emits the PR proof matrix. Findings become a repair queue through \texttt{agent\_fix\_queue}, \texttt{issues export}, or \texttt{repair-plan}. This is the real-time part of Jankurai: not an omniscient daemon, but a repeatable local/PR/CI loop that makes merge evidence current.

\begin{lstlisting}[style=codebox]
jankurai witness . \
  --changed-from origin/main \
  --baseline agent/repo-score.json \
  --out target/jankurai/merge-witness.json \
  --md target/jankurai/merge-witness.md
\end{lstlisting}

The rolling score is the trust ledger behind ratchet mode. \texttt{score diff} compares baseline and candidate reports by score, caps, and finding fingerprint. \texttt{score trend} summarizes recent \texttt{agent/score-history.jsonl} rows so recurring regression is visible before release governance.

The report itself is proof material. JSON and Markdown are the agent repair surface; SARIF, JUnit, GitHub step summaries, JSONL repair queues, and issue exports are integration surfaces. A failed score with no findings is invalid output: every below-floor dimension must produce at least one routed finding with owner, lane, rerun command, evidence kind, fingerprint, and queue item.
